\documentclass[11pt]{article}
\usepackage{amsmath,amssymb,amsthm,mathtools}
\usepackage[margin=1in]{geometry}
\newcommand{\C}{\mathbb C}
\newcommand{\im}{\operatorname{im}}
\newcommand{\rank}{\operatorname{rank}}
\newcommand{\diag}{\operatorname{diag}}
\newcommand{\Span}{\operatorname{span}}
\title{Differentiated original conductors: complete gcds, finite receivers, and leakage ranks}
\author{Independent programme derivation}
\date{20 September 2026}
\begin{document}
\maketitle

\section{Original domain, action and algebraic dual}
This note independently rederives DO19--23 of
\texttt{02\_CONSTANT\_DEPTH\_ORIGINAL\_OBSERVATION.md} and strengthens
their rank bound using the exact repeated-factor profile.  The working
domain is the original simple-quartet period locus stated there: the
conductor has nonzero support on each of its four edges, the original
invariant observation has the stated primitive-column property, and its
five-orbit rank formula is valid.  The original degree is
$k\equiv1\pmod4$, $k\ge77$.  Put $q=(k+1)^2$, $j_-=k-8$, and retain
\[
 \lambda_{ab}=\eta_k+2\delta a+2i\gamma b,\quad
 \eta_k=k(1/2-\delta-i\gamma),\quad
 E_k=\C[S]/\prod_{a,b=0}^k(S-\lambda_{ab}),
 \quad M=M_S.                                                   \tag{DC1}
\]
Here $\delta,\gamma>0$ are the original numbers; in the application
$0<\delta<1/2$, $\gamma>2$.  All roots are distinct because real and
imaginary parts distinguish the two indices.  Thus the value-coordinate
map $\mathsf V_k$ is an isomorphism, with its actual Lagrange inverse.
The observation is the original map $\Lambda_k:E_k\to B_k$ onto its
image, and
\[
 K_k=\ker\Lambda_k,\qquad r=\dim K_k=8k-16.                      \tag{DC2}
\]

Let $R=\C[Z_0,Z_1,W_0,W_1]$, with bidegrees $(1,0)$ for the $Z$
variables and $(0,1)$ for the $W$ variables.  Write $T_{a,b}=R_{a,b}$;
when $a,b\ge0$, its dimension is $(a+1)(b+1)$.  The actual conductor is
\[
 \mathcal A=\sum_{\sigma\in\mathcal S}a_\sigma
 Z_0^{8-\sigma_1}Z_1^{\sigma_1}
 W_0^{8-\sigma_2}W_1^{\sigma_2}\in T_{8,8}.                       \tag{DC3}
\]
Its support and coefficients are unchanged.  For a degree-$(k,k)$
biform $p=\sum p_{ab}Z_0^{k-a}Z_1^aW_0^{k-b}W_1^b$, the relevant
algebraic dual pairing is
\[
 \varphi_p(x)=\sum_{a,b}p_{ab}(\mathsf V_kx)_{ab}.                \tag{DC4}
\]
This is complex linear; it uses a transpose, not a conjugate transpose.
By direct evaluation of DC1, the dual of the actual action $M$ is
\[
 M^\vee=\eta_k+\mathcal D,\qquad
 \mathcal D=2\delta Z_1\partial_{Z_1}
                 +2i\gamma W_1\partial_{W_1}.                   \tag{DC5}
\]
The binomial identities relating $(M^\vee)^n$ and $\mathcal D^n$
retain $\eta_k$ in both directions, so all their row spans through
each fixed order agree.  No original physical action is replaced by a
different one.

Denote by $\mathcal W_k\subset T_{k,k}$ the original invariant
coefficient image, identified with the algebraic dual image of
$\Lambda_k$ by DC4.  The original invariant conductor construction
gives
\[
 \mathcal A T_{k-8,k-8}\subset\mathcal W_k.                       \tag{DC6}
\]
The actual map receiving this inclusion will be made explicit below.

\section{The complete gcd sequence, including every multiplicity}
Factor in the polynomial UFD $R$:
\[
 \mathcal A=c\prod_{\nu=1}^{m}F_\nu^{e_\nu},\quad
 c\in\C^*,\quad \deg F_\nu=(u_\nu,v_\nu),
 \quad e_\nu\ge1.                                               \tag{DC7}
\]
Each factor is bihomogeneous.  To see this directly, highest and
lowest $Z$ degrees in a nonzero product are the sums of the highest
and lowest degrees of its factors, since the polynomial ring is a
domain.  Their differences add and are nonnegative.  A homogeneous
product has difference zero, so each factor is $Z$ homogeneous;
the same argument applies to $W$ degree.  Each nonconstant factor has
$u_\nu+v_\nu\ge1$.  Consequently
\[
 \sum_\nu e_\nu u_\nu=8,\qquad
 \sum_\nu e_\nu v_\nu=8,\qquad m\le16,
 \quad d:=\max_\nu e_\nu\le8.                                   \tag{DC8}
\]
For the last bound, at least one of $u_\nu,v_\nu$ is positive.

Full edge support excludes every coordinate variable as a factor:
the edges $\sigma_1=0,8$ exclude $Z_1,Z_0$, respectively, and the
two other edges exclude $W_1,W_0$.  For any irreducible factor $F$, if
$F\mid\mathcal DF$, equality of bidegrees forces
$\mathcal DF=\zeta F$ for some scalar $\zeta$, including zero.
On a fixed bidegree, the monomial weights are
$2\delta a+2i\gamma b$; they are all distinct.  Thus such an
eigen-polynomial is a single monomial.  An irreducible monomial is
a scalar multiple of one coordinate variable, which is excluded.
We have proved
\[
 F_\nu\nmid\mathcal D F_\nu\quad\hbox{for every }\nu.           \tag{DC9}
\]

Fix $F=F_\nu$, $e=e_\nu$, and write $\mathcal A=F^e U$, $F\nmid U$.
For $0\le l\le e$, the product rule gives the exact leading
$F$-adic term
\[
 \mathcal D^l\mathcal A
 \equiv (e)_{\underline l}F^{e-l}(\mathcal DF)^l U
                          \pmod {F^{e-l+1}}.                    \tag{DC10}
\]
For a proof, at order zero this is immediate.  Differentiating its
leading term produces the stated order-$l+1$ term with coefficient
$(e)_{\underline l}(e-l)$; derivatives of its other factors retain
at least $F^{e-l}$, and differentiating an $F^{e-l+1}$ multiple
retains at least $F^{e-l}$.  This proves the induction through $l=e$.
The coefficient in DC10 is nonzero modulo $F$: the quotient by the
irreducible $F$ is a domain, its characteristic is zero, and both
$\mathcal DF$ and $U$ have nonzero images.  Thus the exact valuation
at order $l\le e$ is $e-l$.

For every $l\ge0$, it follows that the common gcd, up to a nonzero
constant, is exactly
\[
 G_l:=\gcd(\mathcal A,\mathcal D\mathcal A,\ldots,
                       \mathcal D^l\mathcal A)
       =\prod_\nu F_\nu^{\max(e_\nu-l,0)}.                       \tag{DC11}
\]
Indeed every common factor already divides $\mathcal A$, and DC10
computes its smallest valuation among the first $l$ derivatives;
after order $e_\nu$, that smallest valuation is already zero.
In particular $G_d=1$ and $d\le8$, proving DO20 with all repeated
factors retained.  Record its exact bidegree
\[
 g_l=\sum_\nu\max(e_\nu-l,0)u_\nu,\qquad
 h_l=\sum_\nu\max(e_\nu-l,0)v_\nu.                              \tag{DC12}
\]
Equivalently these are the bidegrees of the directly computable gcd
in DC11; no factorization or squarefreeness is assumed in using them.

\section{Finite selection of a pair and its entire biform kernel}
For each $1\le l\le d$, put
\[
 \mathcal B_{l,t}=\sum_{a=1}^{l}t^{a-1}\mathcal D^a\mathcal A.
                                                                    \tag{DC13}
\]
Every summand is divisible by $G_l$.  For every irreducible factor
$F$ of $\mathcal A/G_l$, reduction of
$\mathcal B_{l,t}/G_l$ modulo $F$ is a nonzero polynomial in $t$ of
degree at most $l-1$.  Its nonzero coefficient is supplied by the
valuation calculation DC10 at order $\min(e_F,l)$.
Passing to the fraction field of the domain $R/(F)$ proves that
at most $l-1$ complex values of $t$ make this entire reduction zero.
There are at most $m$ such irreducible factors.  Therefore select the
first integer in
\[
 \{0,1,\ldots,m(l-1)\}
\]
for which the exact gcd test gives
\[
 \gcd(\mathcal A,\mathcal B_l)=G_l.                              \tag{DC14}
\]
The selection exists: the list has one more value than the proved
maximum number of forbidden values.  It retains the actual complex
conductor coefficients.  At $l=8$ the uniform list of 113 integers
in the intake suffices; using $l=d$ requires no higher derivative.
For $l=1$ the selected polynomial is simply $\mathcal D\mathcal A$.

Set $\mathcal F_l=\mathcal A/G_l$ and
$\mathcal H_l=\mathcal B_l/G_l$.  They are coprime, both of
bidegree $(8-g_l,8-h_l)$.  Consider the actual coefficient map
\[
 \mu_l:T_{k-8,k-8}\oplus T_{k-8,k-8}\longrightarrow T_{k,k},
 \qquad (f,g)\longmapsto\mathcal Af+\mathcal B_lg.                \tag{DC15}
\]
Its entire kernel is
\[
 \ker\mu_l=
 \{(\mathcal H_l t,-\mathcal F_l t):
                 t\in T_{k-16+g_l,k-16+h_l}\}.                  \tag{DC16}
\]
For a proof, a kernel element satisfies
$\mathcal F_l f=-\mathcal H_l g$ after cancellation of $G_l$ in
the domain.  Since the two remaining factors are coprime in a UFD,
$\mathcal H_l$ divides $f$ and $\mathcal F_l$ divides $g$.  Substitute
$f=\mathcal H_l t$ and cancel $\mathcal H_l\mathcal F_l$ to obtain
$g=-\mathcal F_l t$.  The bidegrees force exactly the stated degree
of $t$.  The converse follows by multiplication, and the parameter
$t$ is unique.

For $k\ge16$ all displayed degrees are nonnegative.  Hence the
cokernel dimension of the pair map is exactly
\[
 \begin{split}
 c_l&=(k+1)^2-2(k-7)^2
             +(k-15+g_l)(k-15+h_l)\\
    &=128+(g_l+h_l)(k-15)+g_lh_l.                                \end{split}\tag{DC17}
\]
At $l=d$, and in particular at $l=8$, this equals 128.  This is a
complete biform calculation.  Coprimality here does not mean that
the two positive-degree biforms generate the unit ideal in the
four-variable polynomial ring; their finite coefficient cokernel
is precisely what DC17 computes.

\section{The explicit receiver from the original observation stack}
The derivative-row inclusion can be checked without losing the
original output functionals.  If $f\in T_{k-8,k-8}$, then
\[
 (\mathcal D^l\mathcal A)f
  =\sum_{a=0}^l(-1)^{l-a}\binom la
      \mathcal D^a\bigl(\mathcal A\mathcal D^{l-a}f\bigr).
                                                                    \tag{DC18}
\]
To prove it, expand each outer derivative by the product rule.  The
coefficient of $(\mathcal D^b\mathcal A)\mathcal D^{l-b}f$ is
$\binom lb\sum_{a=b}^l(-1)^{l-a}\binom{l-b}{a-b}$, which is one
when $b=l$ and zero otherwise.  DC6 and DC5 then show that the row
space of the actual observation stack through $M^l$ contains
$\im\mu_l$.

Here is the complete map between the original source and target
objects realizing this containment.  Retain $E_{j_-}$ and its
original multiplication $M_-=M_{S'}$.  The finite-shift conductor
is the original map
\[
 C:E_k\to E_{j_-},\qquad C[P]=[T_AP],\qquad
 C=\bar C\Lambda_k.                                             \tag{DC19}
\]
This factorization is well-defined by the original conductor-row
inclusion $\ker\Lambda_k\subset\ker C$; it is unique since
$\Lambda_k$ is onto its image.  In original lower-grid value
coordinates it is exactly
\[
 (Cx)_\beta=\sum_{\sigma\in\mathcal S}
                        a_\sigma(\mathsf V_kx)_{\beta+\sigma},
 \qquad 0\le\beta_1,\beta_2\le k-8.                             \tag{DC20}
\]
Set
\[
 \eta_8=4-8\delta-8i\gamma,
 \quad\omega_\sigma=2\delta\sigma_1+2i\gamma\sigma_2,
 \quad C_0=C,\quad
 C_{a+1}=C_aM-(M_-+\eta_8I)C_a.                                 \tag{DC21}
\]
Every $C_a$ has the same type $E_k\to E_{j_-}$.  The unchanged
upper and lower root values satisfy
\[
 \lambda^{(k)}_{\beta+\sigma}
   -\lambda^{(k-8)}_\beta-\eta_8=\omega_\sigma.
\]
Induction in DC21 therefore proves the exact formula
\[
 (C_ax)_\beta=\sum_\sigma
                  \omega_\sigma^a a_\sigma(\mathsf V_kx)_{\beta+\sigma}.
                                                                    \tag{DC22}
\]
In particular it is the finite-shift conductor whose coefficient
biform is $\mathcal D^a\mathcal A$.  The original constant $\eta_8$
has not been discarded, and none of the original shifts or coefficients
has been replaced by a normalized shift.
Right multiplication by $M$ and left multiplication by
$M_-+\eta_8I$ commute as operators on $\operatorname{Hom}(E_k,E_{j_-})$.
Their binomial expansion gives
\[
 C_a=\sum_{n=0}^a(-1)^{a-n}\binom an
       (M_-+\eta_8I)^{a-n}\bar C\Lambda_k M^n.                   \tag{DC23}
\]

Write $\mathcal O_lx=(\Lambda_kx,\Lambda_kMx,\ldots,
\Lambda_kM^lx)$.  With $t_l$ selected in DC14, define on its actual
output space $B_k^{\oplus(l+1)}$ the map
\[
 \begin{split}
 \mathcal K_l(y_0,\ldots,y_l)=\biggl(&\bar C y_0,\\[-3pt]
 &\sum_{a=1}^lt_l^{a-1}\sum_{n=0}^a(-1)^{a-n}\binom an
            (M_-+\eta_8I)^{a-n}\bar C y_n\biggr)
      \in E_{j_-}\oplus E_{j_-}.
 \end{split}                                                     \tag{DC24}
\]
Then
\[
 \mathcal K_l\mathcal O_l=(C,C_{\mathcal B_l}),
 \qquad C_{\mathcal B_l}=\sum_{a=1}^lt_l^{a-1}C_a.               \tag{DC25}
\]
Pairing the lower-grid values in DC22 with any coefficient biform
$f\in T_{k-8,k-8}$ gives precisely
$\varphi_{(\mathcal D^a\mathcal A)f}$.  Thus the algebraic transpose
of the pair on the right of DC25 is exactly the multiplication map
DC15, under the original value-coordinate isomorphisms.  Its rank
is $q-c_l$, and its kernel dimension is $c_l$.  Therefore
\[
 n_l:=\dim\bigcap_{a=0}^l\ker(\Lambda_kM^a)\le c_l.            \tag{DC26}
\]
This proves DO21 and the stronger multiplicity-resolved bounds.  It
does not assert that the pair rows exhaust the invariant coefficient
image.  All the physical Taylor units, Gamma factors and original
output coordinates reside in the unchanged $\Lambda_k$ and $\bar C$;
DC24 is an actual map on those outputs, not a new observation imposed
by a coefficient analogy.

\section{One-step leakage and the two exact mixed directions}
At the original cutoff $N$, retain the full attained metric $G_N$ on
$E_k$ and the original quotient metric
$Q_N=(\Lambda_kG_N^{-1}\Lambda_k^*)^{-1}$.  Let
\[
 I:\C^r\overset{\sim}{\longrightarrow}K_k,\quad
 H_K=I^*G_NI,\quad L=G_N^{-1}\Lambda_k^*Q_N,
 \quad J_K=H_K^{-1}I^*G_N.
\]
The original decomposition $[I,L]$ is an isometry from
$\diag(H_K,Q_N)$ to $G_N$.  In this frame retain all four blocks
\[
 [I,L]^{-1}M[I,L]=\begin{pmatrix}T&C_{\rm ret}\\B&D\end{pmatrix},
 \quad B=\Lambda_kMI,\quad C_{\rm ret}=J_KML.                    \tag{DC27}
\]
If the initial state is $Ix$, write its two coordinates after $n$
steps as $(x_n,y_n)$.  The full block equations give
\[
 x_{n+1}=Tx_n+C_{\rm ret}y_n,\quad
 y_{n+1}=Bx_n+Dy_n,\quad x_0=x,\quad y_0=0.
\]
Solving the first recursion by induction and substituting it in the
second gives, with every feedback term retained,
\[
 y_{n+1}=BT^nx+Dy_n+
            \sum_{a=0}^{n-1}BT^{n-1-a}C_{\rm ret}y_a.           \tag{DC28}
\]
Thus $(y_1,\ldots,y_l)$ is an invertible lower triangular linear
transform, with identity diagonal, of
$(Bx,BTx,\ldots,BT^{l-1}x)$.  Its inverse is obtained by subtraction
in DC28.  Its kernel dimension is exactly $n_l$: an invisible whole
state lies in $K_k$ already at step zero.  Consequently
\[
 \rank\begin{pmatrix}B\\BT\\\vdots\\BT^{l-1}\end{pmatrix}
   =r-n_l\ge r-c_l,\qquad
 l\,\rank B\ge r-c_l.                                           \tag{DC29}
\]

Define the explicit integer computed from the actual conductor gcds
\[
 R_k^{\rm diff}:=
 \max\left\{0,\max_{1\le l\le8}
 \left\lceil\frac{8k-144-(g_l+h_l)(k-15)-g_lh_l}{l}\right\rceil
                                                    \right\}.
                                                                    \tag{DC30}
\]
Then the original one-step action leakage satisfies
\[
 \boxed{\rank(\Lambda_kMI)\ge R_k^{\rm diff}\ge k-18.}          \tag{DC31}
\]
The second bound uses $g_8=h_8=0$.  The maximum can stop at
$d=\max e_\nu$, since for $l\ge d$ the numerator is the same
positive number $8k-144$ and increasing the denominator cannot
improve the bound.  In particular the exact actual multiplicity
gives the further bound
\[
 \rank B\ge\left\lceil\frac{8k-144}{d}\right\rceil.             \tag{DC32}
\]
For a squarefree original conductor, which is determined by its
actual gcd rather than assumed, $d=1$ and the resulting bound is
$\rank B\ge8k-144$.  Formula DC30 remains valid without any
squarefreeness assertion.  It also retains potential improvements
at intermediate gcd degrees when factors have different multiplicities.

The original full attained-source rank-two Hermitian-weight identity
in DO17 gives, in these exact metrics,
\[
 \widetilde B=Q_N^{1/2}BH_K^{-1/2},\qquad
 \widetilde C=H_K^{1/2}C_{\rm ret}Q_N^{-1/2},\qquad
 \widetilde C+\widetilde B^*=R_N^{\rm wt},\quad
 \rank R_N^{\rm wt}\le2.                                        \tag{DC33}
\]
The invertible metric factors do not change ranks.  The elementary
inequality $\rank(X+Y)\ge\rank X-\rank Y$, proved from
$\im X\subset\im(X+Y)+\im Y$, yields
\[
 \rank C_{\rm ret}\ge\max(0,R_k^{\rm diff}-2).                  \tag{DC34}
\]
For the original metric projection $P_K=IJ_K$, direct block
multiplication gives
\[
 [M,P_K]=\begin{pmatrix}0&-C_{\rm ret}\\B&0\end{pmatrix}.
\]
Its two image summands occupy the disjoint kernel and observation
blocks, so its rank is exactly $\rank B+\rank C_{\rm ret}$.  Hence
\[
 \boxed{\rank[M,P_K]\ge
 R_k^{\rm diff}+\max(0,R_k^{\rm diff}-2)
 =2R_k^{\rm diff}-2\ge2k-38.}                                   \tag{DC35}
\]
The equality of the two displayed bounds uses $k\ge77$, hence
$R_k^{\rm diff}\ge59$.  In particular DO22--23 are correct; DC30,
DC34 and DC35 give their stronger receiving form using the actual
conductor multiplicity profile.

\section{What the identities establish}
The proof supplies an exact finite procedure for the derivative gcds,
an explicit finite selection of the biform pair, its complete syzygy
kernel with both degrees retained, and the actual original observation
receiver DC24.  It evaluates the rank consequences for both mixed
action directions and their commutator.  The first eight steps leave
at most 128 invisible directions; no claim that this remainder is
zero has been made.  The separate DO9 constant-depth reconstruction
removes that remainder using additional original observations.
Neither a rank estimate nor that reconstruction sets the one-step
angle to zero, assigns the complex current signs, or bounds the
smallest nonzero singular value without its original metric calculation.

The direct input examined here is the user's
\texttt{02\_CONSTANT\_DEPTH\_ORIGINAL\_OBSERVATION.md}, DO1--23,
with DO19--23 independently proved above.  The original conductor
factorization in DC19 is also the explicitly proved OB3 map.
The biform, UFD, binomial and finite-matrix arguments used in this
note are given in full; no human result beyond those programme
inputs is invoked without a proof.
\end{document}
